\chapter{Mapa del repositorio}\label{ap:mapa}
\begin{center}
\small
\begin{tabular}{@{}p{0.3\linewidth}p{0.62\linewidth}@{}}
\toprule
\textbf{Ruta} & \textbf{Qué contiene} \\
\midrule
\ruta{server.mjs} & Servidor HTTP, rutas de la API y arranque. \\
\ruta{lib/} & Un módulo por responsabilidad: importador, política de agentes, copias cifradas, correo, PDF, OCR, RAG, lector de webs… \\
\ruta{public/} & Interfaz: \ruta{app.js}, \ruta{core.js}, \ruta{calc.js} (cálculo fiscal), \ruta{commands.js} (órdenes) y \ruta{views/} (una pantalla por archivo). \\
\ruta{test/} & Pruebas con \texttt{node:test}. \\
\ruta{mac/} & Aplicación nativa de macOS en SwiftUI. \\
\ruta{relay/} & Relé ciego y página web del cliente. \\
\ruta{docs/} & Documentación técnica por temas y el tutorial. \\
\ruta{scripts/} & Empaquetado y prueba de carga (\texttt{npm run stress}). \\
\ruta{promo/} & Vídeo de presentación con datos ficticios. \\
\ruta{book/} & Este libro: fuentes en \LaTeX{} y herramientas para fotografiar la aplicación. \\
\bottomrule
\end{tabular}
\end{center}

\chapter{Cómo regenerar el libro}\label{ap:regenerar}
Las capturas se generan contra un entorno con datos ficticios, de forma reproducible:
\begin{codebox}
book/tools/env.sh                       # levanta REM y un relé con datos inventados
node book/tools/shoot-client.mjs        # flujo del cliente (móvil)
node book/tools/shoot.mjs               # pantallas de empresa con marcas numeradas
node book/tools/shoot-mobile.mjs        # pantallas móviles
python3 book/tools/tutorial2tex.py      # Parte III desde docs/tutorial/TUTORIAL.md
book/build.sh                           # compila main.pdf (XeLaTeX)
\end{codebox}
Para añadir una captura nueva, basta con describirla en \ruta{book/tools/shoot.mjs} (ruta, pasos y marcas) y referenciarla desde el capítulo con \verb|\figura|.

% Generado por tools/tutorial2tex.py
\chapter{Glosario}\label{glosario}

\begin{itemize}
\tightlist
\item
  \textbf{ACID}: propiedades de una transacción (atomicidad, consistencia, aislamiento, durabilidad).
\item
  \textbf{AAD (datos asociados)}: datos que se autentican en un cifrado AEAD sin cifrarse.
\item
  \textbf{AES-GCM}: cifrado simétrico con autenticación de integridad.
\item
  \textbf{CSP (Content Security Policy)}: cabecera que limita qué recursos puede cargar una página.
\item
  \textbf{DKIM / SPF}: mecanismos que permiten a un servidor de correo certificar el dominio emisor.
\item
  \textbf{ECDH / HKDF}: intercambio de claves y derivación de una clave de sesión a partir de un secreto compartido.
\item
  \textbf{Long polling}: el cliente pregunta al servidor y el servidor responde cuando hay algo nuevo.
\item
  \textbf{RAG (generación aumentada por recuperación)}: buscar fragmentos relevantes de un corpus y pasarlos al modelo como contexto.
\item
  \textbf{Migración}: cambio de esquema de base de datos numerado y aplicado en orden.
\item
  \textbf{XSS}: inyección de scripts en una página mediante datos no escapados.
\end{itemize}

